% TA-PFC-E1 — Luis Yépez. Matriz de stakeholders del primer avance.
% Fragmento para insertar después de \subsection{Matriz de actores, necesidades y participación}.

La siguiente matriz presenta una propuesta inicial de actores para el sistema de tutorías y seguimiento por indicadores. Los niveles de influencia son estimaciones para planificar la elicitación, no una clasificación aprobada por la institución. Las entrevistas, validaciones y pruebas indicadas son \textbf{actividades previstas}, no realizadas.

\par\medskip
\noindent\begin{minipage}{\textwidth}
\centering
\small
\textbf{Matriz preliminar de partes interesadas del primer avance}\par\smallskip
\begin{tabular}{@{}p{2.5cm}p{3.65cm}p{1.7cm}p{5.3cm}@{}}
\toprule
\textbf{Actor} & \textbf{Necesidad principal} & \textbf{Influencia} & \textbf{Forma de participación prevista} \\
\midrule
Estudiante & Acceder a tutorías y consultar su asistencia, seguimiento e indicadores académicos propios. & Medio & Cuestionarios para conocer necesidades; evaluación de prototipos de solicitud y consulta; pruebas de aceptación de las funciones para estudiantes. \\
\addlinespace
Docente tutor & Organizar sesiones y disponer de información pertinente para identificar dificultades y orientar tutorías. & Alto & Entrevistas sobre el proceso tutorial; revisión de los datos requeridos; validación de registros, criterios de alerta y recomendaciones. \\
\addlinespace
Administrador del sistema & Operar la plataforma de forma segura, con roles y permisos definidos. & Medio & Revisión de restricciones técnicas; evaluación del modelo de accesos; pruebas de administración y seguridad. \\
\addlinespace
Coordinador académico & Supervisar tutorías y contar con indicadores útiles para planificar intervenciones académicas. & Alto & Entrevistas sobre indicadores, reportes y objetivos; priorización y revisión de alertas; participación en aceptación funcional según atribuciones. \\
\bottomrule
\end{tabular}
\end{minipage}
\par\medskip

\paragraph{Justificación de los niveles de influencia.}
Se propone influencia \textbf{alta} para el docente tutor por su conocimiento operativo del acompañamiento académico y para el coordinador por su posible intervención en prioridades y decisiones de seguimiento. Se estima influencia \textbf{media} para el estudiante, por su participación directa en el uso y validación, y para el administrador, por sus decisiones técnicas sin autoridad académica automática. Estos valores deberán modificarse si la institución confirma otra distribución de responsabilidades.

\paragraph{Aspectos pendientes de elicitación.}
Antes de aprobar requisitos debe identificarse quién puede suministrar calificaciones, qué indicadores permiten detectar dificultades sin generar conclusiones engañosas, quién revisa cada recomendación y qué datos puede ver cada perfil. La matriz se limita a los cuatro actores solicitados por el reparto del equipo; las dependencias con fuentes institucionales se documentarán junto con el contexto del sistema que prepara Jordy.
